% Gesamttest «Quadratische Funktionen» — Quelle des PDFs (python3 scripts/build-lp-pdf.py).
% Musterlösung und Raster: bewertungspaket.tex. Alle Werte mit python3 nachgerechnet (02.10.2026).
\documentclass[11pt]{article}
\usepackage{../lp-druck}
\renewcommand{\lpfuss}{Gesamttest · 24 Punkte}
\begin{document}
\lpkopf{Gesamttest}{%
  \vspace{4pt}\noindent\begin{tabularx}{\linewidth}{@{}X X X@{}}
  Code (kein Name): \hrulefill & Datum: \hrulefill & Punkte: \hrulefill\ / 24
  \end{tabularx}}

\begin{lpinfo}{Anleitung}
\textbf{Zeit:} rund 20 Minuten. \textbf{Hilfsmittel:} Teil A und C ohne Taschenrechner, Teil B mit.
Schreib zu jeder Aufgabe den \textbf{Rechenweg} auf — bewertet wird der Weg.
Danach: Lösung fotografieren (ohne Namen) und mit dem \emph{Bewertungspaket} bewerten lassen.
\end{lpinfo}

\lpteil{Teil A · Formen lesen und umformen}{ohne Taschenrechner · 11 Punkte}

\begin{lpaufgabe}{G1}{Scheitelform lesen}{4 P}
$f(x) = -(x-2)^2 + 9$. Gib Scheitelpunkt und Öffnung an, berechne die Nullstellen und schreibe $f$ in der Produktform.
\lpplatz{7}
\end{lpaufgabe}

\begin{lpaufgabe}{G2}{In die Scheitelform}{3 P}
Bringe $f(x) = 2x^2 - 12x + 10$ in die Scheitelform. Gib den Scheitelpunkt und den Schnittpunkt mit der $y$-Achse an.
\lpplatz{6}
\end{lpaufgabe}

\begin{lpaufgabe}{G3}{Vom Graphen zur Gleichung}{2 P}
\begin{minipage}[t]{0.52\linewidth}
Gib die Scheitelform der abgebildeten Parabel an. Die markierten Punkte liegen auf Gitterpunkten.
\lpplatz{4}
\end{minipage}\hfill
\begin{minipage}[t]{0.42\linewidth}\vspace{0pt}\centering
\begin{tikzpicture}
\begin{axis}[width=6.2cm,height=6.2cm,axis lines=middle,xmin=-4,xmax=3,ymin=-3,ymax=4,
  xtick={-4,...,3},ytick={-3,...,4},grid=both,grid style={lplinie},tick label style={font=\scriptsize},
  xlabel=$x$,ylabel=$y$,axis equal image]
\addplot[lpblau,very thick,domain=-3.7:1.6,samples=100] {(x+1)^2-2};
\addplot[only marks,mark=*,color=lpgruen] coordinates {(-1,-2) (0,-1) (1,2)};
\end{axis}
\end{tikzpicture}
\end{minipage}
\end{lpaufgabe}

\begin{lpaufgabe}{G4}{In die Grundform}{2 P}
Schreibe $f(x) = 3(x+1)(x-2)$ in der Grundform $ax^2 + bx + c$.
\lpplatz{3}
\end{lpaufgabe}

\lpteil{Teil B · Aufstellen}{Taschenrechner erlaubt · 6 Punkte}

\begin{lpaufgabe}{G5}{Scheitel und Punkt}{3 P}
Eine Parabel hat den Scheitel $S(3 \mid -2)$ und geht durch $P(5 \mid 6)$. Stelle ihre Funktionsgleichung auf.
\lpplatz{5}
\end{lpaufgabe}

\begin{lpaufgabe}{G6}{Wasserstrahl}{3 P}
Ein Wasserstrahl tritt bei $x = 0$ aus und trifft 8\,m weiter auf dieselbe Höhe; dazwischen steigt er 4\,m hoch.
Stelle die Funktionsgleichung $h(x)$ der Flugbahn auf ($x$ und $h$ in m, Austrittsstelle im Ursprung).
\lpplatz{5}
\end{lpaufgabe}

\lpteil{Teil C · Extremwerte}{ohne Taschenrechner · 7 Punkte}

\begin{lpaufgabe}{G7}{Ballwurf}{3 P}
Ein Ball fliegt nach $h(t) = -5t^2 + 20t + 1.5$ ($t$ in s, $h$ in m). Wann ist er am höchsten, und wie hoch ist er dann?
\lpplatz{5}
\end{lpaufgabe}

\begin{lpaufgabe}{G8}{Regenrinne}{4 P}
Aus einem 24\,cm breiten Blech wird eine Rinne gebogen: links und rechts je $x$\,cm hoch, der Rest ist der Boden.
Bei welchem $x$ ist der rechteckige Querschnitt am grössten? Gib die Masse und die Fläche an.
\lpplatz{7}
\end{lpaufgabe}
\end{document}
